\section*{Вопрос 3}

Основных методов несколько. Контентная фильтрация советует объекты, похожие на то, что пользователь уже выбирал. Коллаборативная фильтрация ищет похожих пользователей и переносит их оценки. Прецедентный подход опирается на типажи и готовые сценарии выбора. Подход на знаниях строится на явных требованиях к объекту и правилах области. Демографическая смотрит на статистику по группам людей. Гибрид смешивает подходы. Сравнение в таблице~\ref{tab:rec_methods}.

\begin{table}[H]
\centering
\caption{Сравнительный анализ методов формирования рекомендаций}
\label{tab:rec_methods}
\scriptsize
\begin{tabular}{|L{2.5cm}|L{2.6cm}|L{2.1cm}|L{2.7cm}|L{2.0cm}|L{2.7cm}|}
\hline
\textbf{Метод} & \textbf{На чём основан} & \textbf{Оценки других пользователей} & \textbf{Холодный старт} & \textbf{Объяснимость} & \textbf{Главная слабость} \\
\hline
Контентная фильтрация & Признаки объектов и профиль пользователя & Не нужны & Новый объект легко. Новому пользователю нужен первый отклик & Высокая. Видно, за какие признаки выдали объект & Предлагает <<ещё того же>> и редко удивляет \\
\hline
Коллаборативная по памяти & Матрица оценок и поиск ближайших соседей & Обязательны & Тяжёлый и для пользователя, и для объекта & Средняя. <<Похожим людям понравилось>> & Разреженные оценки и <<белая ворона>>, чей вкус ни с кем не совпадает \\
\hline
Коллаборативная модельная & Скрытые факторы из матрицы оценок & Обязательны & Тяжёлый для обеих сторон & Низкая. Факторы не имеют смысла & Дорогое обучение и непрозрачный результат \\
\hline
Прецедентная & Типажи пользователей и готовые сценарии выбора & Нужны как основа для типажей & Средний. Типаж работает сразу, но грубо & Средняя. Ссылка на сценарий & Грубая подгонка под типаж, не ловит деталей вкуса \\
\hline
Основанная на знаниях & Явные требования пользователя и правила области & Не нужны & Проблемы нет. Новый объект входит по описанию & Высокая. Ответ следует из требований & Нужна формализация области и диалог с пользователем \\
\hline
Демографическая & Анкетные данные и статистика по группам & Нужны как статистика & Средний. Работает сразу, но грубо & Средняя. Ссылка на категорию людей & Грубые обобщения и вопросы приватности \\
\hline
Гибридная & Комбинация перечисленного & Зависит от состава & Смягчается комбинацией & Зависит от состава & Сложность и настройка весов \\
\hline
\end{tabular}
\end{table}

По проблемам. Холодный старт тяжелее всего у коллаборативных методов и почти отсутствует у контентной и подхода на знаниях. Белая ворона, то есть пользователь, чей вкус ни с кем не совпадает, бьёт только по коллаборативной фильтрации. Контентной она не страшна, потому что профиль строится по признакам объектов, а не по чужим оценкам. Устаревание предпочтений касается всех методов со статичным профилем. Прецедентные системы борются с ней сменой сценария, у остальных нужно отдельное решение.

В моей задаче выбран контентный подход с элементами подхода на знаниях. Пользователь один и матрицы чужих оценок нет, коллаборативным методам нечего считать. Признаки объектов заданы явно, поэтому холодный старт закрыт. Новый ликёр рекомендуется сразу по описанию. Белой вороны нет, профиль строится по контенту. Устаревание гасится тем, что свежие лайки весят больше старых. Фильтры поиска и функция полезности работают как часть подхода на знаниях. Пользователь формулирует требования, а система при пустой выдаче сама советует, что ослабить.

\vspace{1.5em}
\begin{flushright}
Смирнов И.В. ИУ7-11М
\end{flushright}
